\QHeading{第1题　计算填表。}
\par\noindent


\DataTable{4}{交通工具 & 速度/（千米/时） & 时间/时 & 路程/千米\\\hline
汽车 & {\color{answerred}80} & \(12\) & \(960\)\\\hline
火车 & \(112\) & \(16\) & {\color{answerred}1792}\\\hline
摩托车 & \(40\) & {\color{answerred}8} & \(320\)\\\hline}
\par

\par\noindent

\Answer{
\QHeading{计算过程：计算填表}
\begin{center}
\begin{adjustbox}{max width=\linewidth}\begin{tabular}{c|rrr}
\hline
交通工具&速度/（千米/时）&时间/时&路程/千米\\\hline
汽车&80&12&960\\
火车&112&16&1792\\
摩托车&40&8&320\\\hline
\end{tabular}\end{adjustbox}
\end{center}
\begin{enumerate}
\item 汽车速度：\(960\div12=80\)千米/时。
\item 火车路程：\(\begin{aligned}[t]&\phantom{={}}112\times16\\&=112\times(10+6)\\&=112\times10+112\times6\\&=1120+672\\&=1792\text{（千米）}\text{。}\end{aligned}\)
\item 摩托车时间：\(\begin{aligned}[t]&\phantom{={}}320\div40\\&=(320\div10)\div(40\div10)\\&=32\div4\\&=8\text{（小时）}\text{。}\end{aligned}\)
\end{enumerate}
\QHeading{计算过程：第2题②：列车行驶的路程}
路程\(=\)速度\(\times\)时间：
\Formula{\begin{aligned}[t]&\phantom{={}}125\times13\\&=125\times(10+3)\\&=125\times10+125\times3\\&=1250+375\\&=1625\text{（千米）}\text{。}\end{aligned}}
答：甲、乙两地相距1625千米。


}

\QHeading{第2题　解决问题。}

\Problem{①一种超音速飞机2秒可以飞行700米。它的速度是多少？}
\par

\Answer{速度$=$路程$\div$时间。\par $700\div2=350\text{（米/秒）}$。\par 答：它的速度是350米/秒。}
\Problem{②一列火车从甲地开往乙地，平均速度是125千米/时，行驶了13小时到达乙地。甲、乙两地相距多少千米？}

\Problem{③一辆汽车从甲地开往相距255千米的乙地，每小时行驶85千米。这辆汽车一共需要行驶多少小时？}
\Answer{时间$=$路程$\div$速度。\par $255\div85=3\text{（小时）}$。\par 答：这辆车一共需要行驶3小时。}
\Problem{④某运动员骑自行车上山，速度是12千米/时，用了3小时到达山顶。按原路返回用了2小时。返回时的速度是多少？}

\MethodTitle{方法一}
\Answer{\Formula{\begin{aligned}[t]&\phantom{={}}12\times3\div2\\&=36\div2\\&=18\text{（千米/时）}\end{aligned}}\par 答：返回时的速度是18千米/时。}

\MethodTitle{方法二：先求每1小时上山路程折合的下山速度}
\Answer{同一路程原来用3小时，现在用2小时。先平均分配一个小时的12千米，再乘3。}
\Answer{\Formula{\begin{aligned}[t]&\phantom{={}}12\div 2\times 3\\
&=6\times 3\\
&=18\text{（千米/时）}\end{aligned}}}
\Answer{答：返回时的速度是18千米/时。}

